\cleardoublepage
\phantomsection
\begin{planbox}[แผนบริหารการสอนปฏิบัติการที่ 3\quad กฎของนิวตันและการวัดความหนืดตามกฎของสโตกส์ (สัปดาห์ที่ 3)]
\noindent\textbf{1. เวลาเรียนปฏิบัติการ}\quad 3 ชั่วโมง (180 นาที)\quad \textbf{ผลลัพธ์การเรียนรู้ที่ประเมิน}\quad CLO1, CLO2, CLO3, CLO4

\noindent\textbf{2. วัตถุประสงค์เชิงพฤติกรรมประจำปฏิบัติการ}
\begin{enumerate}
    \item อธิบายสมดุลของแรงทั้งสาม (แรงโน้มถ่วง แรงลอยตัว และแรงหนืดของสโตกส์) ที่กระทำต่ออนุภาคทรงกลมขณะจมลงด้วยอัตราเร็วสุดท้ายคงที่ ($v_t$) ได้ \textbf{(CLO1 เชื่อมโยง PLO2)}
    \item ปฏิบัติการวัดอัตราเร็วสุดท้ายของลูกกลมโลหะในของเหลวและคำนวณหาสัมประสิทธิ์ความหนืด ($\eta$) พร้อมเชื่อมโยงสู่การตกตะกอนของเซลล์ยีสต์และแบคทีเรียได้ \textbf{(CLO2, CLO3 เชื่อมโยง PLO1, PLO3)}
\end{enumerate}

\noindent\textbf{3. แผนการจัดเวลาในคาบปฏิบัติการ 3 ชั่วโมง}
\begin{itemize}
    \item \textbf{ชั่วโมงที่ 1}\quad อธิบายสมดุลแรงตามกฎข้อที่ 1 ของนิวตัน ($\sum F = 0$) และสมการกฎของสโตกส์
    \item \textbf{ชั่วโมงที่ 2}\quad วัดรัศมีลูกกลมด้วยไมโครมิเตอร์ และจับเวลาการตกผ่านของเหลวความหนืดต่างกัน (น้ำกลั่น น้ำเชื่อม และกลีเซอรีน)
    \item \textbf{ชั่วโมงที่ 3}\quad คำนวณค่าความหนืด $\eta$ และเปรียบเทียบผลของขนาดรัศมีต่ออัตราเร็วตกตะกอน
\end{itemize}
\end{planbox}
\clearpage

\chapter{กฎของนิวตันและการวัดความหนืดของของเหลวชีวภาพ}
\index{กฎของสโตกส์}\index{สัมประสิทธิ์ความหนืด}\index{อัตราเร็วสุดท้าย}\index{การตกตะกอนของเซลล์}

\section{หลักการทางฟิสิกส์และทฤษฎีเบื้องหลังการทดลอง}

เมื่อปล่อยวัตถุทรงกลมหรือเซลล์จุลินทรีย์รัศมี $r$ ที่มีความหนาแน่น $\rho_p$ ลงในของเหลวที่มีความหนาแน่น $\rho_f$ และมีสัมประสิทธิ์ความหนืด $\eta$ ในช่วงแรกวัตถุจะมีความเร่งลงด้านล่าง แต่เมื่ออัตราเร็วเพิ่มขึ้น แรงต้านความหนืดของของเหลว ($F_{\text{drag}} = 6\pi\eta r v$) จะเพิ่มขึ้นตามจนกระทั่ง\textbf{แรงดึงขึ้นรวม (แรงหนืด $+$ แรงลอยตัว) มีขนาดเท่ากับแรงโน้มถ่วงที่ดึงลงพอดี} ($\sum F = 0$ ตามกฎข้อที่ 1 ของนิวตัน) นับจากจุดนั้นวัตถุจะจมลงด้วย\textbf{อัตราเร็วสุดท้ายคงที่} (\textbf{Terminal Velocity}, $v_t$)

\begin{conceptbox}[สมการคำนวณอัตราเร็วสุดท้ายและสัมประสิทธิ์ความหนืดของสโตกส์]
\begin{enumerate}
    \item \textbf{คำนวณอัตราเร็วสุดท้ายจากการจับเวลา}\quad $v_t = \dfrac{s}{t}\quad (\text{หน่วย m/s})$
    \item \textbf{คำนวณสัมประสิทธิ์ความหนืดของของเหลว} ($\eta$ มีหน่วยเป็น $\text{Pa}\cdot\text{s}$)
    \begin{equation}
        \eta = \frac{2 r^2 (\rho_p - \rho_f) g}{9 v_t}
    \end{equation}
    โดยที่ $g = 9.8\text{ m/s}^2$, $r$ คือรัศมีทรงกลมในหน่วยเมตร ($\text{m}$), $\rho_p$ คือความหนาแน่นทรงกลม ($\text{kg/m}^3$) และ $\rho_f$ คือความหนาแน่นของเหลว ($\text{kg/m}^3$)
\end{enumerate}
\end{conceptbox}

\begin{worked}[การคำนวณความหนืดของกลีเซอรีนจากการจับเวลาลูกเหล็กตก]
\textbf{โจทย์ตัวอย่าง}\quad ลูกเหล็กทรงกลมมีรัศมี $r = 1.5\text{ mm} = 1.5 \times 10^{-3}\text{ m}$ และมีความหนาแน่น $\rho_p = 7{,}800\text{ kg/m}^3$ ถูกปล่อยให้จมลงในกระบอกตวงบรรจุกลีเซอรีน ($\rho_f = 1{,}260\text{ kg/m}^3$) โดยเคลื่อนที่ผ่านระยะห่างระหว่างขีดยาง $s = 0.20\text{ m}$ ในเวลาเฉลี่ย $t = 2.0\text{ s}$ จงหาอัตราเร็วสุดท้าย $v_t$ และความหนืด $\eta$ ของกลีเซอรีน

\textbf{วิธีคำนวณทีละขั้นตอน}
\begin{enumerate}
    \item \textbf{หาอัตราเร็วสุดท้าย}\quad $v_t = \dfrac{0.20\text{ m}}{2.0\text{ s}} = 0.10\text{ m/s}$
    \item \textbf{แทนค่าหาความหนืด $\eta$}
    \[
        \eta = \frac{2 \times (1.5 \times 10^{-3})^2 \times (7{,}800 - 1{,}260) \times 9.8}{9 \times 0.10} \approx 0.32\text{ Pa}\cdot\text{s}
    \]
\end{enumerate}
\end{worked}

\section{แผนภาพชุดอุปกรณ์การทดลองและรายการวัสดุ}

\begin{figure}[H]
\centering
\resizebox{0.82\textwidth}{!}{%
\begin{tikzpicture}[font=\sffamily, >=Stealth]
    % Graduated cylinder
    \fill[boxAmberBg] (2.0,0.4) rectangle (5.2,5.0);
    \draw[physRoyalBlue, line width=1.6pt] (1.6,0.4) -- (5.6,0.4);
    \draw[physRoyalBlue, line width=1.5pt] (2.0,5.2) -- (2.0,0.4) -- (5.2,0.4) -- (5.2,5.2);

    % Rubber bands (Start and Stop markers)
    \draw[crimsonAlert, line width=2.2pt] (1.85,4.0) -- (5.35,4.0) node[right, font=\small\bfseries] {ขีดยางบน (เริ่มกดจับเวลา $t=0$)};
    \draw[crimsonAlert, line width=2.2pt] (1.85,1.4) -- (5.35,1.4) node[right, font=\small\bfseries] {ขีดยางล่าง (หยุดจับเวลาที่ระยะ $s = 0.20\text{ m}$)};
    \draw[<->, physRoyalBlue, line width=1.2pt] (1.3,1.4) -- node[left, font=\small\bfseries] {$s = 20\text{ cm}$} (1.3,4.0);

    % Falling sphere
    \shade[ball color=charcoal!80] (3.6,2.8) circle (0.28cm);
    \draw[->, crimsonAlert, line width=1.4pt] (3.6,2.45) -- (3.6,1.75) node[right, font=\footnotesize\bfseries] {$v_t$ คงที่};
    \draw[->, bioEmerald, line width=1.4pt] (3.6,3.15) -- (3.6,3.85) node[right, font=\footnotesize\bfseries] {$F_{\text{drag}} + F_b$};
\end{tikzpicture}%
}
\caption{การจัดชุดกระบอกตวงวัดความหนืดด้วยวิธีปล่อยลูกกลมโลหะตกผ่านระยะทางมาตรฐาน}
\end{figure}

\begin{equipbox}[รายการอุปกรณ์สำหรับแต่ละกลุ่ม]
\begin{itemize}
    \item กระบอกตวงทรงสูงขนาด $500\text{ mL}$ จำนวน 2 ใบ บรรจุสารละลายน้ำเชื่อมเข้มข้น ($\rho_f = 1{,}200\text{ kg/m}^3$) และกลีเซอรีน ($\rho_f = 1{,}260\text{ kg/m}^3$)
    \item ลูกเหล็กทรงกลมขนาดเล็กและขนาดกลาง ($\rho_p = 7{,}800\text{ kg/m}^3$), ไมโครมิเตอร์, แม่เหล็กสำหรับดึงลูกเหล็กขึ้นจากก้นกระบอก และนาฬิกาจับเวลาดิจิทัล
\end{itemize}
\end{equipbox}

\section{ขั้นตอนการปฏิบัติการทดลอง}
\begin{enumerate}
    \item ใช้ไมโครมิเตอร์วัดเส้นผ่านศูนย์กลาง $D$ ของลูกเหล็กทรงกลม คำนวณรัศมี $r = D/2$ แล้วแปลงหน่วยเป็นเมตร ($\text{m}$ โดยคูณด้วย $10^{-3}$)
    \item ตรวจสอบแถบยางรัดรอบกระบอกตวงให้ขีดบนอยู่ต่ำกว่าผิวของเหลวอย่างน้อย $5\text{ cm}$ และขีดล่างอยู่ห่างจากขีดบนเป็นระยะทาง $s = 20.0\text{ cm}$ ($0.20\text{ m}$)
    \item ใช้ปากคีบปล่อยลูกเหล็กให้จมลงตรงกึ่งกลางกระบอกตวง (ระวังอย่าให้ลูกเหล็กชนผนังแก้ว) กดเริ่มจับเวลาเมื่อลูกเหล็กผ่านขีดยางบน และกดหยุดเมื่อผ่านขีดยางล่าง ทำซ้ำ 3 ครั้งสำหรับของเหลวแต่ละชนิด
\end{enumerate}

\section{ใบบันทึกผลการทดลองประจำกลุ่ม}

\begin{worksheetbox}[ใบบันทึกผลการทดลองที่ 3\quad การวัดความหนืดและอัตราเร็วตกตะกอนตามกฎของสโตกส์]
\noindent\textbf{กลุ่มปฏิบัติการที่} \dotfill \textbf{วันที่ทำการทดลอง} \dotfill\\[4pt]
\noindent\textbf{รัศมีลูกเหล็กที่วัดได้} $r = $ \dotfill $\text{mm} = $ \dotfill $\text{m}$\quad \textbf{ระยะห่างขีดยาง} $s = 0.20\text{ m}$

\vspace{6pt}
\noindent\textbf{ตารางบันทึกผลการจับเวลาการตกและคำนวณสัมประสิทธิ์ความหนืด ($\eta$)}
\begin{center}
\footnotesize\setlength{\tabcolsep}{3.5pt}
\begin{tabularx}{\linewidth}{>{\raggedright\arraybackslash}X c c c c c}
\toprule
\textbf{ชนิดของเหลวชีวภาพ} & \textbf{ครั้งที่ 1 ($\text{s}$)} & \textbf{ครั้งที่ 2 ($\text{s}$)} & \textbf{เฉลี่ย $t$ ($\text{s}$)} & \textbf{$v_t = s/t$ ($\text{m/s}$)} & \textbf{ความหนืด $\eta$ ($\text{Pa}\cdot\text{s}$)} \\
\midrule
สารละลายน้ำเชื่อมเข้มข้น & ............ & ............ & ............ & ............ & \dotfill \\
กลีเซอรีนบริสุทธิ์ & ............ & ............ & ............ & ............ & \dotfill \\
\bottomrule
\end{tabularx}
\end{center}

\vspace{4pt}
\noindent\textbf{แสดงวิธีคำนวณหาค่าความหนืด $\eta$ ของกลีเซอรีนจากข้อมูลที่วัดได้จริง}\\[4pt]
\dotfill\\[6pt]
\dotfill

\noindent\textbf{สรุปและอภิปรายผลการทดลอง}\\[4pt]
\dotfill\\[6pt]
\dotfill
\end{worksheetbox}

\section*{คำถามท้ายปฏิบัติการที่ 3}
\begin{enumerate}
    \item เพราะเหตุใดเราจึงต้องติดขีดยางเส้นบนให้ต่ำลงมาจากผิวหน้าของเหลวประมาณ $5\text{ cm}$ แทนที่จะเริ่มจับเวลาทันทีที่ปล่อยลูกเหล็กที่ผิวของเหลว
    \item ในการเก็บรักษาสายพันธุ์แบคทีเรียในตู้แช่แข็ง $-80\,^\circ\text{C}$ นิยมเติมกลีเซอรีนความเข้มข้น $15\text{--}20\%$ ลงในอาหารเลี้ยงเชื้อ นอกจากป้องกันผลึกน้ำแข็งแล้ว กลีเซอรีนมีผลต่อความหนืดและอัตราการจมตัวของเซลล์อย่างไร
\end{enumerate}
